\section{Problem Statement}
\label{section:31-GHA-Problem-Statement}

There are two driving reasons to develop a benchmark for Graph-Based AI systems. 
First, a user looking to implement a particular algorithm, install a particular accelerator, or adopt a particular programming paradigm must be able to determine which combination of components will best meet their needs~\cite{Carpinetti2002}. 
Second, a developer looking to \textit{exploit} latent capacity in an algorithm, accelerator or programming paradigm or \textit{elevate} the capabilities of one of these systems needs to understand where the bottlenecks exist in the current system~\cite{Saavedra1996, Hooker1994, Hooker1995}.

Towards informing end-user decision-making, a \textit{competitive} benchmark should validate both the \textit{effectiveness} and \textit{efficency} of a system. 
That is, a `high-performing' system should produce provably correct results and describe the amount of resources consumed to produce that result. 
Without balancing these two factors, trivial optimizations to the systems under test can produce incorrect results very quickly, or perfect results extremely inefficently. 

To inform system developers, the results of an \textit{empirical} benchmark must be sufficiently descriptive to illuminate processing bottlenecks. 
Early approaches focused on the \textit{competitive} aspect to the detriment of using the benchmark measurements to guide future development. 
In response to uninformative benchmarks that concealed bottlenecks in amortized program execution costs, the methodology of decomposing programs to their primitive components and measuring each component's performance has become a popular method to expose system bottlenecks~\cite{Saavedra1996}. 
The manifestations of these decomposed programs are as \textit{Atomic Operations}, \textit{Kernels}, \textit{Toy Programs} and \textit{Synthetic Programs}~\cite{Hennessy2019}. 

\textit{Atomic Operation} benchmarking uses very specific operations to develop an understanding of system performance. 
Typically manifesting as \textit{microbenchmarks}, these measure operations like \textit{Floating Point Operations per Second, Traversed Edges Per Second} that provide the building blocks of practical maximal performance that the more comprehensive benchmarks are built on.
For example, a pointer-chasing microbenchmark reveals the practical maximal memory bandwidth of a system~\cite{Beamer2015}.

\textit{Kernel} benchmarking takes the `real' small key pieces of an application and benchmarks their performance by some applicable metric. 
Kernel benchmarks are extremely popular in graph benchmarking, with many systems using BFS, SSSP, PR, CF and CC as the `kernel' operations that they evaluate~\cite{Satish2014, Wang2016, Kang2009, Fu2014, Dadu2021, Beamer2017, Dhulipala2020}.

\textit{Toy Program} benchmarks are the evolution of \textit{Kernel} benchmarks and are typically small programs that one could expect to write in an introduction to programming course~\cite{Hennessy2019}.
They may be composed of several kernels and are representative of a `real' problem, just a very simple one. 
For example, a program to update labels on a graph that needs to load the graph to memory uses the BFS kernel to traverse the graph and update each label in turn before writing it back to storage.

\textit{Synthetic Program} benchmarks are `fake' programs built to simulate the behavior of real programs. 
The general idea is that they are less time and resource intensive to construct as a mock-up, and so can be used to evaluate a system  `realistically' without needing to completely build the application under test. 